A massive matter unit is a composite of two conformal spinors (twistors) whose flat mass is the bracket, $m^{2} = 2|\br{\pi^{1}\pi^{2}}|^{2}$ at $\Lambda = 0$. The bracket transforms by $\det(U)$ under the internal $U(2)$, and under the global form with the hypercharge of Fact \claimref{neutrino.imported-derived-content-singlets} the bracket of two conformal spinors, each the vacuum of its pair system, transforms as $(1,1)$ with $Y = +1$: the quantum numbers of $e^{c}$, and not those of a weak doublet. The constituents there are slot-occupying states (that paper's $(1,0)$ components, which exist in its state space under Hypothesis (S), each of hypercharge $+\tfrac12$, so that the $+1$ is their sum; the slot arithmetic $\lambda^{3}$ per slot is unconditional in the global form), not the slot-free grade-zero state, the empty conformal spinor (this paper's empty twistor); the bracket is defined on such conformal spinors, is computed there only for that case, is not defined for an excited composite, and says nothing, either way, about a bilinear of two grade-zero states. The relation between this object and the Standard Model's order parameter is not settled in the series, and the reading of the bracket as an order parameter is recorded at observation level only. \emph{Provenance:} \cite[Props.~3.1-3.2]{ButlerT12} (\claimref{standard-model.mass-bracket-charged-singlet}, \claimref{standard-model.internal-symmetry-ladder}); \cite[Prop.~4.1]{ButlerT11} (\claimref{mass.mass-symplectic-area-two-twistor-formula}); the form of this Fact is the families paper's \cite[Fact~2.1]{ButlerT13} (\claimref{families.imported-face-valued-bracket-quantum-numbers}).
